% --- Archon physics formalization source begin ---
% archon:physics
% archon:covers PhyXMiniProblems/problem_phyx_mini_0199.lean
% archon:source-report reports/phyx_mini/problem_phyx_mini_0199.source.json

\chapter{Physics problem phyx\_mini\_0199}
\label{ch:PhyXMiniProblems_problem_phyx_mini_0199}

\paragraph{Problem source.}
\#\# Physical scenario

When a family of four with a total mass of \textbackslash{}(200\textbackslash{},\textbackslash{}mathrm\{kg\}\textbackslash{}) step into their \textbackslash{}(1200\textbackslash{},\textbackslash{}mathrm\{kg\}\textbackslash{}) car, the car's springs compress \textbackslash{}(3.0\textbackslash{},\textbackslash{}mathrm\{cm\}\textbackslash{}). The car was loaded with \textbackslash{}(300\textbackslash{},\textbackslash{}mathrm\{kg\}\textbackslash{}) rather than \textbackslash{}(200\textbackslash{},\textbackslash{}mathrm\{kg\}\textbackslash{}).

\#\# Question

How far will the car lower?

\#\# Answer choices

- A: A: \$7.5\textbackslash{}times10\textasciicircum{}4\textbackslash{}

- B: B: \$6.5\textbackslash{}times10\textasciicircum{}3\textbackslash{}

- C: C: \$5.5\textbackslash{}times10\textasciicircum{}4\textbackslash{}

- D: D: \$6.5\textbackslash{}times10\textasciicircum{}4\textbackslash{}

\#\# Figure caption (auxiliary; use the image as primary evidence)

The image depicts a detailed view of a MacPherson strut suspension system, commonly used in automotive vehicles. Key components visible include:

1. **Coil Spring**: Positioned at the top, aiding in absorbing shock.
2. **Shock Absorber/Strut**: Beneath the coil spring, painted blue, responsible for dampening vibrations.
3. **Control Arm**: Curved metallic arm connecting various parts, allowing the wheel to move up and down independently.
4. **Brake Disc/Rotors**: Circular component at the bottom, part of the braking system.
5. **Brake Caliper**: Attached to the brake disc, containing brake pads to create friction for stopping.
6. **Ball Joint**: At the end of the control arm, allowing pivoting movement.
7. **Mounting Bolts**: Seen throughout the structure, securing parts together.

The suspension system is crucial for ensuring stability and comfort by managing wheel movement and absorbing impacts.

\#\# Dataset metadata

Wave Properties; Physical Model Grounding Reasoning, Predictive Reasoning

\paragraph{Recorded answer/context.}
D: D: \$6.5\textbackslash{}times10\textasciicircum{}4\textbackslash{}

\paragraph{Figure/image path.}
/root/proposal\_for\_physic/science-mango/phyx\_mini\_run/phyx\_data/test\_image/199.png

\paragraph{Formalization target.}
create a compiling Lean file with sorry bodies at `PhyXMiniProblems/problem\_phyx\_mini\_0199.lean`. The Lean declarations must preserve the physical quantities, dimensions or dimensional roles, figure labels, governing-law hypotheses, and final relation expressed by this problem.
Use Mathlib/Physlib names found through LeanExplore where available. If a physics API is missing, introduce faithful local abstractions rather than scalar placeholder aliases.

\begin{theorem}[Physics formalization target]
\label{thm:physics:phyx_mini_0199:target}
\lean{PhyXMiniProblems.ProblemPhyXMini0199.problem_phyx_mini_0199}
\uses{def:physics:phyx-mini-0199:phyxminiproblems-problemphyxmini0199-lengthincentimeters, def:physics:phyx-mini-0199:phyxminiproblems-problemphyxmini0199-carsuspensionsetup, def:physics:phyx-mini-0199:phyxminiproblems-problemphyxmini0199-matchesproblemstatement, def:physics:phyx-mini-0199:phyxminiproblems-problemphyxmini0199-matchessuppliedmacphersonfigure, def:physics:phyx-mini-0199:phyxminiproblems-problemphyxmini0199-hasphysicalsuspensionparameters, def:physics:phyx-mini-0199:phyxminiproblems-problemphyxmini0199-satisfiesstaticsuspensionlaws}
The assigned autoformalize agent should translate this physics problem into Lean declarations in the covered file, with theorem and lemma proof bodies written as `by sorry`.
\end{theorem}
\begin{proof}
This is an autoformalization task, not a proof task. Produce statements that can later be proved by the physics prover without weakening the physical model.
\end{proof}

% --- Archon declaration topology begin ---
\section{Declaration topology}
\begin{definition}[Mass Quantity]
  \label{def:physics:phyx-mini-0199:phyxminiproblems-problemphyxmini0199-massquantity}
  \lean{PhyXMiniProblems.ProblemPhyXMini0199.MassQuantity}
  A physical mass represented coherently in every choice of units
\end{definition}
\begin{proof}
  This is the declared model definition of Mass Quantity.
\end{proof}
\begin{definition}[Length Quantity]
  \label{def:physics:phyx-mini-0199:phyxminiproblems-problemphyxmini0199-lengthquantity}
  \lean{PhyXMiniProblems.ProblemPhyXMini0199.LengthQuantity}
  A physical vertical displacement or compression
\end{definition}
\begin{proof}
  This is the declared model definition of Length Quantity.
\end{proof}
\begin{definition}[Acceleration Quantity]
  \label{def:physics:phyx-mini-0199:phyxminiproblems-problemphyxmini0199-accelerationquantity}
  \lean{PhyXMiniProblems.ProblemPhyXMini0199.AccelerationQuantity}
  A physical acceleration, used for the local gravitational acceleration
\end{definition}
\begin{proof}
  This is the declared model definition of Acceleration Quantity.
\end{proof}
\begin{definition}[Spring Stiffness Quantity]
  \label{def:physics:phyx-mini-0199:phyxminiproblems-problemphyxmini0199-springstiffnessquantity}
  \lean{PhyXMiniProblems.ProblemPhyXMini0199.SpringStiffnessQuantity}
  Equivalent stiffness of the car's parallel suspension springs
\end{definition}
\begin{proof}
  This is the declared model definition of Spring Stiffness Quantity.
\end{proof}
\begin{definition}[mass Readout]
  \label{def:physics:phyx-mini-0199:phyxminiproblems-problemphyxmini0199-massreadout}
  \lean{PhyXMiniProblems.ProblemPhyXMini0199.massReadout}
  \uses{def:physics:phyx-mini-0199:phyxminiproblems-problemphyxmini0199-massquantity}
  Scalar mass readout in a selected coherent unit system
\end{definition}
\begin{proof}
  This is the declared model definition of mass Readout.
\end{proof}
\begin{definition}[length Readout]
  \label{def:physics:phyx-mini-0199:phyxminiproblems-problemphyxmini0199-lengthreadout}
  \lean{PhyXMiniProblems.ProblemPhyXMini0199.lengthReadout}
  \uses{def:physics:phyx-mini-0199:phyxminiproblems-problemphyxmini0199-lengthquantity}
  Scalar length readout in a selected coherent unit system
\end{definition}
\begin{proof}
  This is the declared model definition of length Readout.
\end{proof}
\begin{definition}[acceleration Readout]
  \label{def:physics:phyx-mini-0199:phyxminiproblems-problemphyxmini0199-accelerationreadout}
  \lean{PhyXMiniProblems.ProblemPhyXMini0199.accelerationReadout}
  \uses{def:physics:phyx-mini-0199:phyxminiproblems-problemphyxmini0199-accelerationquantity}
  Scalar acceleration readout in a selected coherent unit system
\end{definition}
\begin{proof}
  This is the declared model definition of acceleration Readout.
\end{proof}
\begin{definition}[spring Stiffness Readout]
  \label{def:physics:phyx-mini-0199:phyxminiproblems-problemphyxmini0199-springstiffnessreadout}
  \lean{PhyXMiniProblems.ProblemPhyXMini0199.springStiffnessReadout}
  \uses{def:physics:phyx-mini-0199:phyxminiproblems-problemphyxmini0199-springstiffnessquantity}
  Scalar spring-stiffness readout in a selected coherent unit system
\end{definition}
\begin{proof}
  This is the declared model definition of spring Stiffness Readout.
\end{proof}
\begin{definition}[mass In Kilograms]
  \label{def:physics:phyx-mini-0199:phyxminiproblems-problemphyxmini0199-massinkilograms}
  \lean{PhyXMiniProblems.ProblemPhyXMini0199.massInKilograms}
  \uses{def:physics:phyx-mini-0199:phyxminiproblems-problemphyxmini0199-massquantity, def:physics:phyx-mini-0199:phyxminiproblems-problemphyxmini0199-massreadout}
  The kilogram readout supplied by the SI unit choice
\end{definition}
\begin{proof}
  This is the declared model definition of mass In Kilograms.
\end{proof}
\begin{definition}[centimeter Unit Choices]
  \label{def:physics:phyx-mini-0199:phyxminiproblems-problemphyxmini0199-centimeterunitchoices}
  \lean{PhyXMiniProblems.ProblemPhyXMini0199.centimeterUnitChoices}
  SI base units with centimeters selected as the length unit
\end{definition}
\begin{proof}
  This is the declared model definition of centimeter Unit Choices.
\end{proof}
\begin{definition}[length In Centimeters]
  \label{def:physics:phyx-mini-0199:phyxminiproblems-problemphyxmini0199-lengthincentimeters}
  \lean{PhyXMiniProblems.ProblemPhyXMini0199.lengthInCentimeters}
  \uses{def:physics:phyx-mini-0199:phyxminiproblems-problemphyxmini0199-lengthquantity, def:physics:phyx-mini-0199:phyxminiproblems-problemphyxmini0199-lengthreadout, def:physics:phyx-mini-0199:phyxminiproblems-problemphyxmini0199-centimeterunitchoices}
  The centimeter readout used for the measured and requested lowering
\end{definition}
\begin{proof}
  This is the declared model definition of length In Centimeters.
\end{proof}
\begin{definition}[Suspension Component]
  \label{def:physics:phyx-mini-0199:phyxminiproblems-problemphyxmini0199-suspensioncomponent}
  \lean{PhyXMiniProblems.ProblemPhyXMini0199.SuspensionComponent}
  Mechanical components visible in the supplied suspension bitmap
\end{definition}
\begin{proof}
  This is the declared model definition of Suspension Component.
\end{proof}
\begin{definition}[Suspension Role]
  \label{def:physics:phyx-mini-0199:phyxminiproblems-problemphyxmini0199-suspensionrole}
  \lean{PhyXMiniProblems.ProblemPhyXMini0199.SuspensionRole}
  Qualitative physical roles attributed to the visible components
\end{definition}
\begin{proof}
  This is the declared model definition of Suspension Role.
\end{proof}
\begin{definition}[Suspension Kind]
  \label{def:physics:phyx-mini-0199:phyxminiproblems-problemphyxmini0199-suspensionkind}
  \lean{PhyXMiniProblems.ProblemPhyXMini0199.SuspensionKind}
  Suspension architecture identified by the primary image
\end{definition}
\begin{proof}
  This is the declared model definition of Suspension Kind.
\end{proof}
\begin{definition}[Compression Reference]
  \label{def:physics:phyx-mini-0199:phyxminiproblems-problemphyxmini0199-compressionreference}
  \lean{PhyXMiniProblems.ProblemPhyXMini0199.CompressionReference}
  The reference height from which the stated extra lowering is measured
\end{definition}
\begin{proof}
  This is the declared model definition of Compression Reference.
\end{proof}
\begin{definition}[Suspension Figure]
  \label{def:physics:phyx-mini-0199:phyxminiproblems-problemphyxmini0199-suspensionfigure}
  \lean{PhyXMiniProblems.ProblemPhyXMini0199.SuspensionFigure}
  \uses{def:physics:phyx-mini-0199:phyxminiproblems-problemphyxmini0199-suspensioncomponent, def:physics:phyx-mini-0199:phyxminiproblems-problemphyxmini0199-suspensionrole}
  Qualitative visibility, geometry, and component roles read from the image. The bitmap contains no mass or compression scale
\end{definition}
\begin{proof}
  This is the declared model definition of Suspension Figure.
\end{proof}
\begin{definition}[Car Suspension Setup]
  \label{def:physics:phyx-mini-0199:phyxminiproblems-problemphyxmini0199-carsuspensionsetup}
  \lean{PhyXMiniProblems.ProblemPhyXMini0199.CarSuspensionSetup}
  \uses{def:physics:phyx-mini-0199:phyxminiproblems-problemphyxmini0199-massquantity, def:physics:phyx-mini-0199:phyxminiproblems-problemphyxmini0199-lengthquantity, def:physics:phyx-mini-0199:phyxminiproblems-problemphyxmini0199-accelerationquantity, def:physics:phyx-mini-0199:phyxminiproblems-problemphyxmini0199-springstiffnessquantity, def:physics:phyx-mini-0199:phyxminiproblems-problemphyxmini0199-suspensionkind, def:physics:phyx-mini-0199:phyxminiproblems-problemphyxmini0199-compressionreference, def:physics:phyx-mini-0199:phyxminiproblems-problemphyxmini0199-suspensionfigure}
  The car, the two added-load cases, and the common suspension model. additionalLowering is measured relative to the equilibrium height of the car without either added load. It is an independent response function; its value at requestedAddedMass is not fixed by this structure
\end{definition}
\begin{proof}
  This is the declared model definition of Car Suspension Setup.
\end{proof}
\begin{definition}[Matches Problem Statement]
  \label{def:physics:phyx-mini-0199:phyxminiproblems-problemphyxmini0199-matchesproblemstatement}
  \lean{PhyXMiniProblems.ProblemPhyXMini0199.MatchesProblemStatement}
  \uses{def:physics:phyx-mini-0199:phyxminiproblems-problemphyxmini0199-massinkilograms, def:physics:phyx-mini-0199:phyxminiproblems-problemphyxmini0199-lengthincentimeters, def:physics:phyx-mini-0199:phyxminiproblems-problemphyxmini0199-carsuspensionsetup}
  Numerical data stated in the prose. In particular, only the 200 kg calibration compression is recorded; the compression under 300 kg is absent
\end{definition}
\begin{proof}
  This is the declared model definition of Matches Problem Statement.
\end{proof}
\begin{definition}[Matches Supplied Mac Pherson Figure]
  \label{def:physics:phyx-mini-0199:phyxminiproblems-problemphyxmini0199-matchessuppliedmacphersonfigure}
  \lean{PhyXMiniProblems.ProblemPhyXMini0199.MatchesSuppliedMacPhersonFigure}
  \uses{def:physics:phyx-mini-0199:phyxminiproblems-problemphyxmini0199-carsuspensionsetup}
  Readouts from the supplied MacPherson-strut image. These fields preserve the visible components and their qualitative geometry, but assert no numerical lowering and therefore do not determine the answer
\end{definition}
\begin{proof}
  This is the declared model definition of Matches Supplied Mac Pherson Figure.
\end{proof}
\begin{definition}[Has Physical Suspension Parameters]
  \label{def:physics:phyx-mini-0199:phyxminiproblems-problemphyxmini0199-hasphysicalsuspensionparameters}
  \lean{PhyXMiniProblems.ProblemPhyXMini0199.HasPhysicalSuspensionParameters}
  \uses{def:physics:phyx-mini-0199:phyxminiproblems-problemphyxmini0199-accelerationreadout, def:physics:phyx-mini-0199:phyxminiproblems-problemphyxmini0199-springstiffnessreadout, def:physics:phyx-mini-0199:phyxminiproblems-problemphyxmini0199-massinkilograms, def:physics:phyx-mini-0199:phyxminiproblems-problemphyxmini0199-lengthincentimeters, def:physics:phyx-mini-0199:phyxminiproblems-problemphyxmini0199-carsuspensionsetup}
  Positivity and nondegeneracy conditions for the static loading regime
\end{definition}
\begin{proof}
  This is the declared model definition of Has Physical Suspension Parameters.
\end{proof}
\begin{definition}[Satisfies Static Suspension Laws]
  \label{def:physics:phyx-mini-0199:phyxminiproblems-problemphyxmini0199-satisfiesstaticsuspensionlaws}
  \lean{PhyXMiniProblems.ProblemPhyXMini0199.SatisfiesStaticSuspensionLaws}
  \uses{def:physics:phyx-mini-0199:phyxminiproblems-problemphyxmini0199-massquantity, def:physics:phyx-mini-0199:phyxminiproblems-problemphyxmini0199-massreadout, def:physics:phyx-mini-0199:phyxminiproblems-problemphyxmini0199-lengthreadout, def:physics:phyx-mini-0199:phyxminiproblems-problemphyxmini0199-accelerationreadout, def:physics:phyx-mini-0199:phyxminiproblems-problemphyxmini0199-springstiffnessreadout, def:physics:phyx-mini-0199:phyxminiproblems-problemphyxmini0199-massinkilograms, def:physics:phyx-mini-0199:phyxminiproblems-problemphyxmini0199-carsuspensionsetup}
  Static Hooke equilibrium for the *additional* load on the unchanged car: k\_eff x(m) = m g. It is stated in every coherent unit system and for every nonnegative added mass, so the single equivalent-stiffness field applies to both load cases. This governing law contains neither 4.5 cm nor any value for the requested compression
\end{definition}
\begin{proof}
  This is the declared model definition of Satisfies Static Suspension Laws.
\end{proof}
\begin{lemma}[added Mass compression cross relation]
  \label{lem:physics:phyx-mini-0199:phyxminiproblems-problemphyxmini0199-addedmass-compression-cross-relation}
  \lean{PhyXMiniProblems.ProblemPhyXMini0199.addedMass_compression_cross_relation}
  \uses{def:physics:phyx-mini-0199:phyxminiproblems-problemphyxmini0199-massreadout, def:physics:phyx-mini-0199:phyxminiproblems-problemphyxmini0199-lengthreadout, def:physics:phyx-mini-0199:phyxminiproblems-problemphyxmini0199-carsuspensionsetup, def:physics:phyx-mini-0199:phyxminiproblems-problemphyxmini0199-hasphysicalsuspensionparameters, def:physics:phyx-mini-0199:phyxminiproblems-problemphyxmini0199-satisfiesstaticsuspensionlaws}
  Eliminating the common stiffness and gravitational acceleration gives the cross-multiplied proportionality between the calibration and requested load cases. No numerical value for the requested lowering is assumed
\end{lemma}
\begin{proof}
  The conclusion follows from the governing relations and figure readouts named in the statement of added Mass compression cross relation.
\end{proof}
\begin{definition}[Answer Choice]
  \label{def:physics:phyx-mini-0199:phyxminiproblems-problemphyxmini0199-answerchoice}
  \lean{PhyXMiniProblems.ProblemPhyXMini0199.AnswerChoice}
  Labels attached to the four answer tokens in the source record
\end{definition}
\begin{proof}
  This is the declared model definition of Answer Choice.
\end{proof}
\begin{definition}[displayed Raw Numeric Token]
  \label{def:physics:phyx-mini-0199:phyxminiproblems-problemphyxmini0199-displayedrawnumerictoken}
  \lean{PhyXMiniProblems.ProblemPhyXMini0199.displayedRawNumericToken}
  \uses{def:physics:phyx-mini-0199:phyxminiproblems-problemphyxmini0199-answerchoice}
  The raw numeric tokens printed by the source. Since no units accompany them and their magnitudes do not represent the requested centimeter lowering, they are retained only as dimensionless dataset metadata
\end{definition}
\begin{proof}
  This is the declared model definition of displayed Raw Numeric Token.
\end{proof}
\begin{definition}[recorded Answer Choice]
  \label{def:physics:phyx-mini-0199:phyxminiproblems-problemphyxmini0199-recordedanswerchoice}
  \lean{PhyXMiniProblems.ProblemPhyXMini0199.recordedAnswerChoice}
  \uses{def:physics:phyx-mini-0199:phyxminiproblems-problemphyxmini0199-answerchoice}
  The answer label recorded in the source dataset
\end{definition}
\begin{proof}
  This is the declared model definition of recorded Answer Choice.
\end{proof}
% --- Archon declaration topology end ---
% --- Archon physics formalization source end ---
